\chapter{Preface} \label{ch:preface}
\index{Preface@\emph{Preface}}

%\section{What this book is about?}\label{sec:aboutThisBook}

This book comprises organic programming experiences, home-grown development techniques,
and other observed patterns, many of which are not mentioned elsewhere.
This book is also a summary of the struggles and the learning, the stumbles and the gains, the dark
quests and the `aha' moments recorded in my software engineering journey over the years.

\section{Who am I?}\label{sec:whoAmI}

When my doctoral research came to a complete stall after six hard-working years, I was convinced
that I had no future as an unachieved physicist.
Digging into my old dreams, I was delighted to find that my passion for coding,
dormant throughout the years, is still unquenchable.


Through self-directed, incremental, and prolonged learning, I built up my software development
skills, one at a time.
It took me a decade or so to grow from a mere passionate dilettante to an experienced professional.
During this period, I read vast amount of books, articles, blog posts on software engineering
best practices, design patterns, programming paradigms, \texttt{etc.}
I implemented a number of fundamental algorithms and data structures.
I attempted to solve as many coding problems as I can find.
And I put them to work, which I deem as the ultimate touchstone for learning.


But more important than these, which every ambitious software developer would do,
I have an insatiable passion for one other thing---I periodically reflect on my
past code and review my notes.
I would do anything to get to the bottom of things that bothers me.
I follow my `guts feeling'---I would not give in if `something in there' upsets me.


Many of the problems that I worked on were revisited up to five or more times over the years.
Each time I was able to gain new understandings and each time my solution shrinks.
I try different solutions to solve the same problem again and again and
I keep asking myself ``Why my implementation works or does not work?''


I painstakingly dedicate quality time to think over bugs I made, bugs I fixed,
and bugs I strategically avoided.
I contemplate on new insights and attempt to reconcile it with previous experience.
I draw analogies across seemingly different fields.


Have you ever wondered why comparison-based sorting algorithms are unanimously bound
by $O(N \log N)$?
I managed to gain a psychologically satisfying understanding through enduring contemplations.
Starting with \texttt{a} sorting algorithm,
I related sorting operations as pumping Shannon's information entropy out of the subject of sorting,
similar to pumping heat out of the room with AC\cite{Shannon-mathematical-1948}.
From there further up, I track it back to the field of statistical
mechanics, developed by Ludwig Boltzmann in the 19th century.
Finally I settled my quest with modeling the sorting time complexity by the
statistical entropy of an $N$-particle isolated system\cite{landau:1980}.
By the way, this result is both surprising and expected to me.


My journey of learning and practicing in the field of computer sciences was materialized in
thousands of Evernote\textregistered pages, $40,000$
lines of coding exercise on problems found online \footnote{\textit{e.g.},
\url{projecteuler.net}, \url{leetcode.com}, \url{glassdoor.com}} or offline,
and more than $100,000$ lines of production code for companies that I worked for.


\section{Why this book?
Why now?}\label{sec:whyThisBook}

I thought about writing a book like this long ago.
But I didn't take this seriously until fairly recently when
I came across a web article talking about the then latest research report by Forrester
which claims that the prolonged shortage in software engineering talent has turned from
`lack of quantity' to `lack of quality'
\footnote{\url{https://go.forrester.com/research/predictions}}.

Quoting the article directly\footnote{Titled as ``2018's Software Engineering Talent
Shortage---It's quality, not just quantity'' on \url{https://hackernoon.com}}:
\begin{quote}
    ``The software engineering shortage is not a lack of individuals calling
    themselves `engineers', the shortage is one of quality---a lack of well-studied,
    experienced engineers with a formal and deep understanding of software engineering.''
\end{quote}
That report resonated with my own observation.
Starting about five years ago,
I started to note that many other developers, colleagues of mine included,
struggled with the same problems that I used to struggle with.
I've worked with graduates and post-graduates from top-tier schools
through numerous design reviews, code reviews, and technical interviews.
They have eye-catching GPAs.
Their resumes are full of buzz words.
They spent a great deal of time reading books on software development best practices, principles,
patterns, guidance, and advice.
But by end of the day when these best practices are to be put to production code, they
needed help.

As a result, it is not surprising to me that codebases, proprietary and open source alike,
all that I have ever worked on without exception,
are filled with copy-n-pasted code that could have been avoided with a little shared function,
intricate logic that could have been simplified into a one-liner,
sneaky bugs caused by an escaped variable from its due scope,
mammoth solutions that could have been refactored into a group of lean, single-purposed functions,
and so on.
It is like glimpsing out of the window of the airplane in a flight, you see a duct-taped wing.
Browsing through implementations of fairly dense algorithms, it is not rare to see
the `smells' listed in Martin Fowler's seminal book `Refactor'~\cite{Fowler1999}.
This may be shocking or even outrageous to mention but true that your mission-critical transactions
actually run on smelly code.
Nowadays ugly code has become the norm rather than exception.
As a matter of fact, if I see any code close to those in `Programming
Pearls' and its sort\cite{Bentley:1986:PP}, I would be surprised.

These plain facts convey the crying signal that there is a gap between the domain of best practices
and the domain how to practice the best practices.
Good news is that the body of `best practices' already exists.
Bad news is that putting the `best practices' to work itself takes years of \emph{practice}.
Let's face the reality---it is a daunting task to attaining the `best practices' in ones own code.
The real-world problems is usually much more complex than the toy examples used in
`best practices' books.
Unless they have first-hand empirical evidence,
developers would not \emph{feel} the values and limits of many of the `best practices'.
Unless they get their hands dirty with fitting multiple principles, design patterns and other
`best practice' tips on a couple of real-world projects, they would not understand the trade-offs,
the collaborations, and the conflicts among various `best practices'.


\section{Who should read this book?} \label{sec:whoshouldread}

Therefore, the sole focus for this book is to serve as an exemplary work on ``How to practice the
best practices in software development?''
By reproducing the battles I fought and sharing the insights I gained, I hope to
bridge the body of software development best practices with developers
on the industrial battle field.

I hope readers would find this book to be somewhat expected as well somewhat unexpected.
As Jon Bentley put it ``This sure looks familiar \ldots but I've never seen this
before''\cite{Bentley:1986:PP}.
This book is intended for professionals who have already had a solid command of
computer science nonetheless who still find application of the software development best practices
in kaleidoscopic coding situations challenging.
It is assumed that readers are comfortable working with one or two of the popular programming
languages such as \python, \java, C and/or C++.
Basic knowledge in algorithm, data structure, and operating system would be a plus but is not
critical to comprehend the key ideas.


\section{What this book is NOT about?}\label{sec:thisBookIsNotAbout}

Some disclaimers are needed here to disperse untrue assumptions for this book.
This book is \emph{not} a programming primer at least not systematically.
Nor does it teach you a specific programming language.


This book is not about recipe or cookbook for any specific framework or system.
Rather, it's a general guideline for development with all frameworks and systems.


Also this book does not offer problem-solving strategies or
quick-n-dirty programming `tips' designed for coding interviews.
Nevertheless, if absorbed into one's coding habit, the techniques in this book \emph{do} help
with build up one's code sense and instinct which, in turn, helps with technical interviews
in the long run.


Last but not least, this book is not a typical preparation material for coding interviews.
Coding interview emphasizes primarily on `Get it done'.
But the focus of this book is to `Get it done right'.
Also many of the techniques presented in this book such as ~\fullref{ch:organizeYourIfElse} or
~\fullref{ch:reusableFunctions} requires more time than allowed in an interview setting.
So don't be bound by them in an interview setting.


Nevertheless, anyone who are striving for long-term investment of
improving their programming skills may find this book inspiring and beneficial, including
those who are preparing for technical interviews.
Coding is a life-long endeavor and it takes both creativity and perseverance.
This book shall be taken as seed rather than food, inspiration rather than formula.


\section{Acknowledgement}

A lot of people to thank\ldots


\begin{flushright}
    S.~W. @ San Mateo College Library
\end{flushright}
